\section{Step 59 Exercise boundaries and equivalent schedules}\label{h263-quant-s09}
Run the public feedback against your working source:
\begin{verbatim}
python3 construction/h263-quantization-construction/grade.py \
    my-quantizer.asm --cli /path/to/mars-cli.jar \
    --report quantizer-results.json
\end{verbatim}
The public cases test three quantizer parameters, both block types, modes and
byte orders, several input patterns, invalid fields, domain failures and state
preservation. Input begins at an address aligned for a halfword but not a vector,
so a direct vector load from the caller needs its own alignment reasoning.
Staging levels privately avoids that input alignment issue in the taught path.

Try another iteration order or private frame. Results and access boundaries
must stay the same. A different schedule is acceptable; the final exercise
contract does not require the reference's labels or scratch registers.

\paragraph{Observe} Successful feedback checks all 128 output halfwords. It
also verifies that a late invalid coefficient leaves the complete result unchanged.
\Needspace{8\baselineskip}
\paragraph{Check} What does exact kernel output agreement establish?
\begin{itemize}
\item[A.] The whole stream is already conforming H.263.
\item[B.] Every complete encoder must choose these same forward levels.
\item[C.] The tested kernel follows this exercise's quantization policy.
\end{itemize}
